%=============================================================================
%  02_introduction.tex  -  Chapter 1: Introduction
%=============================================================================
\chapter{Introduction}
\label{ch:intro}

\section{Motivation}
The last few years have made it trivially easy to \emph{ask a language model a
question}, and correspondingly easy for that model to answer with fluent,
confident, and \emph{wrong} numbers. For data analysis this failure mode is
particularly dangerous: a plausible-looking figure or a neat chart carries an
authority that a paragraph of prose does not. A tool that ``chats with your
data'' is only useful if the analyst can trust that what it reports is actually
\emph{in} the data.

Universal Analyst is built around that trust problem. Rather than treating the
language model as an oracle, it treats it as a \emph{reasoning engine wrapped in
guardrails}: the data is profiled deterministically, relevant facts are
retrieved from a vector index built from the user's own file, and every chart
the model proposes is validated against the real schema before it is rendered.
The model plans; the system verifies; the interface never fakes.

\begin{uanote}[The core promise]
Given a dataset and a question, the analyst gets a correct, grounded answer or a
real chart --- \emph{fast} --- and can trust that nothing was fabricated. Empty
states and error states say what actually happened.
\end{uanote}

\section{Problem Statement}
Existing ``AI + data'' tools tend to fall into one of three traps:

\begin{enumerate}
  \item \textbf{Hallucinated analysis.} The model invents statistics or trends
        that are not supported by the data, because it never actually reads the
        data --- only a truncated preview or nothing at all.
  \item \textbf{Cloud dependence.} The dataset is shipped to a third-party API,
        raising privacy, cost, and reproducibility concerns that are
        unacceptable for many analysts and organizations.
  \item \textbf{Decorative dashboards.} Charts are generated from templates that
        look impressive but are disconnected from the underlying columns, or are
        silently substituted when the model fails.
\end{enumerate}

Universal Analyst addresses all three: it grounds answers in retrieval over the
uploaded data, it runs the model entirely \emph{locally} on the user's GPU
runtime, and it renders only charts that survive a strict server-side
validation pass.

\section{Objectives}
\begin{table}[h]
\centering
\caption{Primary project objectives and how they are met.}
\begin{tabularx}{\linewidth}{@{}L{4.7cm} X@{}}
\toprule
\textbf{Objective} & \textbf{Realization in the system} \\
\midrule
Grounded question answering & RAG over a per-session ChromaDB index built from
schema, per-column, and row chunks of the uploaded file. \\
Local, private inference & Gemma\,4 (12B primary, \code{e4b} router) served via
Ollama on the same GPU box; no data leaves the runtime. \\
Trustworthy visualization & The model emits a strict chart \emph{plan}; the
\code{ChartFactory} validates columns and types before rendering Plotly. \\
Multi-format ingestion & A single \code{IngestionEngine} handling CSV, JSON,
Parquet, PDF tables, and SQLite. \\
Actionable outputs & Skills for cleaning, prediction, and a branded PDF report,
in addition to chat and dashboards. \\
Operational honesty & The \emph{no-fabrication rule}: real streamed output or a
real error, never a placeholder. \\
\bottomrule
\end{tabularx}
\end{table}

\section{Scope}
This project delivers a complete, runnable full-stack application: a Python
\code{3.12} FastAPI backend, a React\,18 single-page front end, a retrieval
layer, a local-LLM integration, a validated visualization engine, a prediction
pipeline, a PDF reporting engine, container definitions for development and
production, a one-click cloud-setup notebook, and a continuous-integration
workflow. It targets a single-analyst, single-session analytical workflow on a
cloud GPU runtime (for example, Lightning.ai). Multi-tenant scaling, persistent
user accounts, and a managed database are explicitly out of scope for this
iteration and are discussed as future work in Chapter~\ref{ch:future}.

\section{Contributions}
\begin{itemize}
  \item A \textbf{grounded analytical assistant} that combines deterministic
        profiling with retrieval-augmented generation over the user's data.
  \item A \textbf{validated, model-driven visualization pipeline} that turns a
        constrained LLM chart plan into a real, correct Plotly figure.
  \item A \textbf{skill-routing architecture} that lets one natural-language
        interface dispatch to analysis, visualization, cleaning, prediction, and
        reporting.
  \item A \textbf{fully local, streaming inference path} with token-level usage
        accounting and multimodal (image) support.
  \item A \textbf{reproducible deployment story} spanning a setup notebook,
        Docker Compose (development and production), and CI.
\end{itemize}

\section{Document Roadmap}
Chapter~\ref{ch:product} frames the product philosophy and users.
Chapter~\ref{ch:arch} presents the end-to-end architecture and request
lifecycle. Chapters~\ref{ch:ingestion}--\ref{ch:reporting} descend into each
backend subsystem in the order data flows through it: ingestion and profiling,
the RAG layer, the LLM client, the skill system, visualization, prediction, and
reporting. Chapter~\ref{ch:frontend} covers the React front end;
Chapter~\ref{ch:api} the HTTP surface; Chapter~\ref{ch:deploy} configuration and
deployment; Chapter~\ref{ch:testing} testing and quality;
Chapter~\ref{ch:security} security and limitations. The document closes with the
team split (Chapter~\ref{ch:team}), future work (Chapter~\ref{ch:future}), a
conclusion, and appendices.
